\documentclass[conference]{IEEEtran}
\IEEEoverridecommandlockouts

\usepackage{cite}
\usepackage{amsmath,amssymb,amsfonts}
\usepackage{algorithmic}
\usepackage{graphicx}
\usepackage{textcomp}
\usepackage{xcolor}
\usepackage{booktabs}
\usepackage{url}
\usepackage{listings}
\usepackage{algorithm}

\lstset{
  basicstyle=\ttfamily\scriptsize,
  breaklines=true,
  frame=single,
  columns=fullflexible,
  captionpos=b
}

\def\BibTeX{{\rm B\kern-.05em{\sc i\kern-.025em b}\kern-.08em
    T\kern-.1667em\lower.7ex\hbox{E}\kern-.125emX}}

\begin{document}

\title{Summit: A Multiplayer AI Coding Agent with Persistent Project Memory\\
\thanks{This research was supported by the Department of Computer Science and Engineering.}
}

\author{\IEEEauthorblockN{Muhilan S\textsuperscript{1}, Dhanesh S K\textsuperscript{2}}
\IEEEauthorblockA{\textit{Department of Computer Science and Engineering} \\
\textit{Rajalakshmi Engineering College}\\
Chennai, India \\
\textsuperscript{1}muhilan.s.2024.csd@rajalakshmi.edu.in, \textsuperscript{2}dhanesh.sk.2024.csd@rajalakshmi.edu.in}
}

\maketitle

\begin{abstract}
Contemporary AI coding assistants operate in a stateless, single-user paradigm. Each interactive session starts in isolation without awareness of prior architectural decisions, team conventions, or project evolution. Furthermore, existing tools lack support for shared, real-time AI collaboration among multiple human developers interacting on the same codebase simultaneously. This paper presents \textbf{Summit}, a collaborative multiplayer AI coding workspace designed to address both limitations through two core mechanisms: (1) a persistent project memory system backed by an ACID-compliant structured SQLite database that retains architectural decisions, coding standards, and execution history across sessions; and (2) a multiplayer WebSocket room architecture that synchronizes multiple human developers with a shared AI coding agent. Summit incorporates an asynchronous Python agent engine built on FastAPI, a React/TypeScript/Monaco workspace frontend, codebase-aware context retrieval, and optimistic concurrency control for concurrent file edits. The agent executes a multi-turn tool-calling loop (supporting workspace file I/O, terminal execution, structured memory updates, and Git checkpoints) routed through LiteLLM. Forensic source code analysis and a 17-test functional verification suite validate system correctness. Comprehensive engineering trade-offs, formal algorithms, and proposed empirical benchmarks are presented.
\end{abstract}

\begin{IEEEkeywords}
AI coding agent, persistent project memory, multiplayer collaboration, WebSocket, LLM tool calling, optimistic concurrency control, context retrieval, FastAPI.
\end{IEEEkeywords}

\section{Introduction}
Large Language Model (LLM)-based coding assistants, such as GitHub Copilot \cite{copilot}, Amazon CodeWhisperer \cite{codewhisperer}, and OpenAI ChatGPT \cite{chatgpt}, have redefined software engineering workflows. Despite significant individual productivity enhancements \cite{peng2023copilot}, state-of-the-art tools exhibit two fundamental structural deficiencies:

\begin{enumerate}
    \item \textbf{The Episodic-Temporal Deficit:} Conventional assistants operate ephemerally. Prior architectural choices, convention agreements, or debugging discoveries are lost across disjoint sessions \cite{hatalis2024memory}. Developers are forced into repetitive prompt engineering to re-explain domain context.
    \item \textbf{The Single-User Isolation Paradigm:} Software engineering is fundamentally collaborative \cite{ko2006exploratory}. While teams interact concurrently over shared repositories, AI coding assistants remain private single-tenant tools. Existing platforms provide no mechanism for distributed developers to co-interact with a shared agent or inspect real-time agent telemetry.
\end{enumerate}

To bridge this gap, we present \textbf{Summit}, an open, extensible multiplayer AI coding workspace with persistent project memory. Summit enables distributed human developers to connect concurrently to a collaborative project room and interact with a unified AI coding agent. The agent maintains long-term project awareness via a dual-tier SQLite memory architecture, pairing conversation history with structured project knowledge. An automated context assembly pipeline extracts project memories, conversation logs, and path-relevance-ranked codebase files into each prompt.

\subsection{Key Contributions}
The main contributions of this paper include:
\begin{itemize}
    \item \textbf{Multiplayer Agent Room Architecture:} A real-time FastAPI WebSocket engine that broadcasts all agent actions (tool calls, file mutations, terminal streams, lifecycle state transitions) to all room participants concurrently.
    \item \textbf{Dual-Memory Persistence Engine:} An ACID-compliant relational memory model decoupling ephemeral conversational logs (\texttt{messages}) from persistent architectural facts (\texttt{project\_memories}) with upsert semantics.
    \item \textbf{Codebase Context Retrieval Algorithm:} A deterministic path-relevance algorithm that injects relevant files, recent conversation, and the full project memory corpus into LLM prompts without expensive embedding overhead.
    \item \textbf{Optimistic Concurrency Control:} Integer-based file versioning detecting edit collisions between developers and agents, returning standard HTTP 409 responses with conflict resolution payloads.
    \item \textbf{Forensic Analysis and Formalization:} Complete algorithmic definitions for context assembly and agent execution loops, validated against 17 integration test cases.
\end{itemize}

\section{Related Work}

\subsection{AI Coding Assistants and Autonomous Agents}
Modern AI developer tools have progressed from line-level completion to autonomous multi-step agents. GitHub Copilot \cite{copilot} and CodeWhisperer \cite{codewhisperer} rely strictly on active in-editor buffers. SWE-agent \cite{yang2024sweagent} introduced Agent-Computer Interfaces (ACI) for resolving real-world GitHub issues \cite{jimenez2024swebench}. Devin \cite{devin2024} and OpenDevin \cite{wang2024opendevin} demonstrated end-to-end task execution in sandboxed shells, while CodeAct \cite{wang2024codeact} pioneered executable Python actions. However, these systems operate in single-tenant, non-persistent environments.

\subsection{Memory in LLM Architectures}
Zhong et al. \cite{zhang2024memorysurvey} categorize agent memory into working, episodic, semantic, and procedural components. Standard RAG architectures \cite{lewis2020rag, shrivastava2023rlcg, zhang2023repocoder} retrieve snippets using dense vector embeddings over static repositories. In contrast, Summit adopts structured, human-interpretable relational memory authored dynamically by both developers and agents, ensuring full injection determinism.

\subsection{Collaborative Development Environments}
Real-time collaborative editing is typically governed by Operational Transformation (OT) \cite{ellis1989ot} or Conflict-Free Replicated Data Types (CRDTs) \cite{oster2006crdt}. While systems like VS Code Live Share operate at the keystroke layer, Summit implements file-level optimistic concurrency control tailored to human-agent co-authoring.

\begin{figure}[t]
\centering
\fbox{
\begin{minipage}{0.92\linewidth}
\scriptsize
\textbf{Client Tier (React 19 / TypeScript / Monaco Editor / Xterm)}
\begin{center}
$\updownarrow$ \textit{HTTP REST (CRUD \& Git) + WebSocket (Real-Time Streams)}
\end{center}
\textbf{FastAPI Application Server (Uvicorn ASGI)}
\begin{itemize}
    \item \textbf{API Routing Layer:} 9 Routers (Projects, Files, Agent, Git, WS, Auth)
    \item \textbf{Agent Engine (SummitAdapter):} LiteLLM Multi-Provider Dispatcher
    \item \textbf{SessionManager:} In-Memory WS Room Registry \& Lifecycle Events
    \item \textbf{ContextRetriever:} Path-Relevance Scoring \& Prompt Assembler
    \item \textbf{WorkspaceManager:} Sandboxed Filesystem Access \& Conflict Check
    \item \textbf{GitService:} Subprocess-Driven Local Checkpoints \& Remote Sync
\end{itemize}
\begin{center}
$\updownarrow$ \textit{Direct SQLite Connection Pool (Zero ORM Overhead)}
\end{center}
\textbf{Relational Persistence Tier (SQLite: summit.db)}
\begin{itemize}
    \item \texttt{users}, \texttt{projects}, \texttt{messages}, \texttt{project\_memories}, \texttt{file\_versions}, \texttt{git\_checkpoints}
\end{itemize}
\end{minipage}
}
\caption{Summit Full-Stack Multiplayer Architecture.}
\label{fig:arch}
\end{figure}

\section{System Design and Methodology}

\subsection{Architecture Overview}
As depicted in Fig.~\ref{fig:arch}, Summit is designed as a decoupled, multi-tier system. The backend is implemented in asynchronous Python (FastAPI/Uvicorn), and the client is a single-page React 19 application featuring Monaco Editor, interactive memory drawers, and an integrated terminal.

\subsection{Multiplayer WebSocket Room Synchronization}
Project-level concurrency is governed by the \texttt{SessionManager}. Active connections are stored in an in-memory dictionary:
\begin{equation}
\mathcal{C}: \text{ProjectID} \rightarrow \{ \text{WS}_i \mapsto (\text{UserID}_i, \text{DisplayName}_i) \}
\end{equation}

When a client establishes a connection to \texttt{/api/projects/\{id\}/agent/stream}, the server:
\begin{enumerate}
    \item Accepts the WebSocket handshake.
    \item Registers the socket in $\mathcal{C}(\text{ProjectID})$ and creates or retrieves the user profile.
    \item Broadcasts a \texttt{USER\_JOINED} event to all connected peers and returns the current \texttt{AGENT\_STATUS}.
\end{enumerate}
Substantive messages follow a server-to-client broadcast paradigm. To prevent race conditions, only one active LLM task can execute per project room at any instant ($\text{Lock}(\text{ProjectID})$).

\subsection{Dual-Tier Relational Memory Formulation}
Summit establishes two distinct persistence vectors within SQLite:
\begin{itemize}
    \item \textbf{Conversation Memory ($\mathcal{M}_{\text{conv}}$):} An append-only sequence of user instructions and agent responses $\mathcal{M}_{\text{conv}} = \langle m_1, m_2, \dots, m_T \rangle$.
    \item \textbf{Project Memory ($\mathcal{M}_{\text{proj}}$):} A set of structured key-value triples:
    \begin{equation}
    \mathcal{M}_{\text{proj}} = \{ (c_k, k, v_k) \mid k \in \mathcal{K} \}
    \end{equation}
    where $c_k \in \{\text{architecture}, \text{decision}, \text{convention}, \text{task}\}$, enforce by a unique constraint $\text{UNIQUE}(project\_id, key)$ executing upsert operations:
    \begin{equation}
    \text{INSERT} \dots \text{ON CONFLICT}(project\_id, key) \text{ DO UPDATE SET } v = v_{\text{new}}
    \end{equation}
\end{itemize}

\subsection{Optimistic Concurrency Control (OCC)}
To prevent conflicting overwrites between human developers and autonomous tool-writes, every file record maintains an integer version $V \in \mathbb{N}^+$. A client update request $R = (p, C_{\text{new}}, V_{\text{exp}})$ succeeds if and only if $V_{\text{exp}} = V_{\text{curr}}$. Upon mismatch, the server returns HTTP 409 (\texttt{Conflict}):
\begin{equation}
\Delta = \text{Diff}(C_{\text{server}}, C_{\text{client}})
\end{equation}
enabling the client to either adopt server state or force-overwrite.

\begin{algorithm}[t]
\caption{Codebase Context Assembly}
\label{alg:context}
\begin{algorithmic}[1]
\REQUIRE $p_{\text{id}}$: Project ID, $q$: User Prompt
\ENSURE $\mathcal{X}$: Structured Prompt String
\STATE $\mathcal{M}_{\text{proj}} \leftarrow \text{Repository.get\_memories}(p_{\text{id}})$
\STATE $\mathcal{M}_{\text{conv}} \leftarrow \text{Repository.get\_messages}(p_{\text{id}}, \text{limit}=10)[-6:]$
\STATE $\mathcal{F}_{\text{all}} \leftarrow \text{Workspace.list\_files}(p_{\text{id}})$
\STATE $\mathcal{T} \leftarrow \{ t \in \text{lower}(q).\text{split}() \mid |t| > 2 \}$
\FORALL{$f \in \mathcal{F}_{\text{all}}$}
    \STATE $S(f) \leftarrow 0$
    \IF{$f \text{ matches standard entrypoints (e.g., main.py, README)}$}
        \STATE $S(f) \leftarrow S(f) + 2$
    \ENDIF
    \FORALL{$t \in \mathcal{T}$}
        \IF{$t \text{ is substring of } \text{lower}(f)$}
            \STATE $S(f) \leftarrow S(f) + 5$
        \ENDIF
    \ENDFOR
\ENDFOR
\STATE $\mathcal{F}_{\text{sel}} \leftarrow \text{TopRanked}(\mathcal{F}_{\text{all}}, S, K=4)$
\STATE $\mathcal{C}_{\text{files}} \leftarrow []$
\FORALL{$f \in \mathcal{F}_{\text{sel}}$}
    \STATE $\text{lines} \leftarrow \text{ReadFile}(f).\text{splitlines}()$
    \STATE $\text{content} \leftarrow \text{Join}(\text{lines}[:200]) + (\text{len}(\text{lines}) > 200 \mathrel{?} \text{ ``[Truncated]'' : ``''})$
    \STATE $\mathcal{C}_{\text{files}}.\text{append}(f \oplus \text{content})$
\ENDFOR
\STATE $\mathcal{X} \leftarrow \text{Format}(\mathcal{M}_{\text{proj}}, \mathcal{M}_{\text{conv}}, \mathcal{C}_{\text{files}}, q)$
\RETURN $\mathcal{X}$
\end{algorithmic}
\end{algorithm}

\section{Implementation Details}

\subsection{Agent Tool Loop and Dispatcher}
The autonomous loop (Algorithm~\ref{alg:context}) runs asynchronously within an \texttt{asyncio.Task}. Tool definitions exposed to the model include:
\begin{enumerate}
    \item \texttt{read\_file(path)}: Retrieves workspace file contents.
    \item \texttt{write\_file(path, content)}: Persists file and broadcasts \texttt{FILE\_CHANGED}.
    \item \texttt{run\_terminal(command)}: Spawns sandboxed subprocesses, capturing stdout/stderr.
    \item \texttt{save\_memory(key, value, category)}: Persists architectural facts to SQLite.
    \item \texttt{git\_checkpoint(message)}: Creates point-in-time local Git snapshots.
\end{enumerate}

\subsection{Database Schema}
The persistence layer comprises six tables:
\begin{itemize}
    \item \texttt{users} (\underline{id}, display\_name, created\_at)
    \item \texttt{projects} (\underline{id}, name, template, workspace\_path, updated\_at)
    \item \texttt{messages} (\underline{id}, project\_id, user\_id, role, content, timestamp)
    \item \texttt{project\_memories} (\underline{id}, project\_id, category, key, value, updated\_at)
    \item \texttt{file\_versions} (\underline{project\_id, path}, version, content, updated\_by)
    \item \texttt{git\_checkpoints} (\underline{id}, project\_id, commit\_hash, message, created\_at)
\end{itemize}

\begin{table}[t]
\caption{System Subsystems and Claim Traceability Matrix}
\label{tab:traceability}
\centering
\scriptsize
\begin{tabular}{lllp{2.5cm}}
\toprule
\textbf{Subsystem} & \textbf{Module File} & \textbf{Component / Function} & \textbf{Status} \\
\midrule
Multiplayer WS Room & \texttt{api/ws.py} & \texttt{agent\_event\_stream} & Verified Passing \\
Session State & \texttt{agent/session\_manager.py} & \texttt{SessionManager} & Verified Passing \\
Conversation Store & \texttt{database/repository.py} & \texttt{save\_message} & Verified Passing \\
Project Memory & \texttt{database/repository.py} & \texttt{save\_memory (upsert)} & Verified Passing \\
Context Assembler & \texttt{context/retriever.py} & \texttt{build\_context} & Verified Passing \\
Concurrency Engine & \texttt{workspace/manager.py} & \texttt{write\_file (OCC)} & Verified Passing \\
Tool Dispatcher & \texttt{agent/summit\_adapter.py} & \texttt{\_execute\_real\_llm\_agent} & Verified Passing \\
Git Checkpointing & \texttt{git/git\_service.py} & \texttt{create\_checkpoint} & Verified Passing \\
Path Sandboxing & \texttt{workspace/manager.py} & \texttt{validate\_safe\_path} & Verified Passing \\
\bottomrule
\end{tabular}
\end{table}

\section{Verification and Benchmark Methodology}

\subsection{Functional Test Suite}
The system implementation was verified using a 17-test functional and integration test suite across four test suites:
\begin{itemize}
    \item \texttt{test\_api.py} (4 tests): Verifies project CRUD lifecycles, file endpoints, and agent health APIs.
    \item \texttt{test\_event\_normalizer.py} (4 tests): Validates Pydantic schema serialization across 15 WebSocket event types.
    \item \texttt{test\_multiplayer\_memory.py} (5 tests): Confirms SQLite memory upserts, OCC version collision rejection (HTTP 409), and Git checkpoint integrity.
    \item \texttt{test\_workspace.py} (4 tests): Enforces workspace path traversal isolation (e.g., blocking \texttt{../../etc/passwd}) and directory tree listings.
\end{itemize}

\subsection{Proposed Empirical Benchmark Protocol}
To quantitatively evaluate collaborative multi-agent workspaces, we formulate the following standardized benchmark metrics:
\begin{enumerate}
    \item \textbf{Memory Recall Fidelity ($R_{\text{mem}}$):} Fraction of previously recorded architectural constraints correctly satisfied in subsequent multi-turn sessions:
    \begin{equation}
    R_{\text{mem}} = \frac{|\mathcal{K}_{\text{adhered}}|}{|\mathcal{K}_{\text{stored}}|}
    \end{equation}
    \item \textbf{Collaboration Collision Rate ($C_{\text{occ}}$):} Frequency of edit conflicts generated during asynchronous human-agent editing under synthetic load.
    \item \textbf{Multiplayer End-to-End Latency ($\tau_{95}$):} 95th percentile event broadcast delay across distributed WebSocket clients.
\end{enumerate}

\section{Limitations and Future Work}
While Summit provides a fully operational multiplayer AI programming platform, several engineering trade-offs warrant future enhancement:
\begin{enumerate}
    \item \textbf{Semantic Memory Pruning:} Summit currently injects all stored memories into prompt sections without token budgeting. Integrating vector embedding stores (e.g., ChromaDB) with semantic similarity filtering will prevent context window overflow.
    \item \textbf{Keystroke-Level CRDT Editing:} Upgrading file-level OCC to fine-grained Conflict-Free Replicated Data Types (CRDTs) will enable simultaneous in-line typing between developers and the agent.
    \item \textbf{Sandboxed Code Execution:} Incorporating microVMs or Docker containers will protect host operating systems from untrusted agent shell commands.
\end{enumerate}

\section{Conclusion}
This paper presented Summit, a multiplayer AI coding agent workspace with persistent project memory. By combining real-time WebSocket room multiplexing, dual-tier SQLite memory persistence, deterministic context retrieval, and optimistic concurrency control, Summit overcomes the isolation and statelessness of current AI developer tools. Forensic source verification and functional test suites confirm architectural viability, laying the groundwork for collaborative human-AI team engineering.

\begin{thebibliography}{00}

\bibitem{copilot} GitHub, ``GitHub Copilot: Your AI pair programmer,'' GitHub Inc., 2021. [Online]. Available: \url{https://copilot.github.com}

\bibitem{codewhisperer} Amazon Web Services, ``Amazon CodeWhisperer AI coding companion,'' AWS, 2022.

\bibitem{chatgpt} OpenAI, ``Introducing ChatGPT,'' OpenAI, 2022. [Online]. Available: \url{https://openai.com/blog/chatgpt}

\bibitem{hatalis2024memory} M. Hatalis et al., ``Memory matters: The need to improve long-term memory in LLM-agents,'' in \textit{Proc. AAAI Workshop on Cognitive Architectures for LLM Agents}, 2024.

\bibitem{ko2006exploratory} A. Ko, B. Myers, M. Coblenz, and H. Aung, ``An exploratory study of how developers seek, relate, and collect relevant information during software maintenance tasks,'' \textit{IEEE Trans. Software Eng.}, vol. 32, no. 12, pp. 971--987, Dec. 2006.

\bibitem{peng2023copilot} S. Peng, E. Kalliamvakou, P. Croft, and M. Counts, ``The impact of AI on developer productivity: Evidence from GitHub Copilot,'' \textit{arXiv preprint arXiv:2302.06590}, 2023.

\bibitem{yang2024sweagent} J. Yang, C. Jimenez, A. Wettig, K. Lieret, S. Yao, K. Narasimhan, and O. Press, ``SWE-agent: Agent-computer interfaces enable automated software engineering,'' \textit{arXiv preprint arXiv:2405.15793}, 2024.

\bibitem{jimenez2024swebench} C. Jimenez et al., ``SWE-bench: Can language models resolve real-world GitHub issues?'' in \textit{Proc. ICLR}, 2024.

\bibitem{devin2024} Cognition AI, ``Introducing Devin, the first AI software engineer,'' Cognition AI, 2024.

\bibitem{wang2024opendevin} X. Wang et al., ``OpenDevin: An open platform for AI software developers as generalist agents,'' \textit{arXiv preprint arXiv:2407.16741}, 2024.

\bibitem{wang2024codeact} X. Wang et al., ``Executable code actions elicit better LLM agents,'' \textit{arXiv preprint arXiv:2402.01030}, 2024.

\bibitem{zhang2024memorysurvey} Z. Zhang, B. Chen, F. Liu, and B. Liu, ``A survey on the memory mechanism of large language model based agents,'' \textit{arXiv preprint arXiv:2404.13501}, 2024.

\bibitem{lewis2020rag} P. Lewis et al., ``Retrieval-augmented generation for knowledge-intensive NLP tasks,'' in \textit{Proc. NeurIPS}, 2020, pp. 9459--9474.

\bibitem{shrivastava2023rlcg} D. Shrivastava, H. Larochelle, and D. Tarlow, ``Repository-level prompt generation for large language models of code,'' in \textit{Proc. ICML}, 2023.

\bibitem{zhang2023repocoder} F. Zhang et al., ``RepoCoder: Repository-level code completion through iterative retrieval and generation,'' in \textit{Proc. EMNLP}, 2023.

\bibitem{ellis1989ot} C. Ellis and S. Gibbs, ``Concurrency control in groupware systems,'' in \textit{Proc. ACM SIGMOD}, 1989, pp. 399--407.

\bibitem{oster2006crdt} G. Oster, P. Urso, P. Molli, and A. Imine, ``Data consistency for P2P collaborative editing,'' in \textit{Proc. ACM CSCW}, 2006, pp. 259--267.

\bibitem{fastapi} S. Ramirez, ``FastAPI: High performance, easy to learn, fast to code, ready for production,'' 2018. [Online]. Available: \url{https://fastapi.tiangolo.com}

\end{thebibliography}

\end{document}
